\documentclass[10pt,a4paper]{article}

\usepackage{amssymb}
\usepackage{amsmath}

\begin{document}

\section{Introduction}

\begin{itemize}
\item \textbf{Slide 1} \par
\noindent Merci M. le Pr\'esident, \\
\noindent Merci Messieurs et Dames, membres du jury. \par
\noindent Bonjour et bienvenue \`a tous.
\noindent Je suis ONGONWOU Fred\'eric, physicien, et aujourd'hui
je vais vous entretenir sur mon sujet de th\`ese 
qui a trait \`a la physique atomique.
Sp\'ecifiquement, mon sujet de th\`ese porte sur la dynamique d'une particule 
dans un champ laser intense.
Nous proposons une \textbf{nouvelle approche th\'eotique 
dans l'espace des moments}. 
Ce travail a \'et\'e ex\'ecut\'e sous la supervision
du Pr. EKOGO Thierry Blanchard du Laboratoire d'Optique et Laser
de l'Universit\'e des Sciences des Techniques de Masuku
et membre du jury.  

\item \textbf{Slide 2} \par
Notre travail s'inscrit dans la cadre plus g\'en\'eral de 
ce qui est connu la "Physique Atomique et Mol\'eculaire et Optique
Non-Lin\'eaire". 
Cette branche de la physique est d'une importance capitale pour 
le d\'eveloppements des sciences et des technologies 
dans la mesure o\`u elle a une incidence directe sur presque toutes 
les avanc\'ees techniques depuis l'invention du laser dans les ann\'ees 
1960s. Tout indique \'egalement que l'optique non lin\'eaire restera 
au centre du d\'eveloppement pour les prochaines d\'ecennies \`a venir.
En effet, on peut voir que la recherche se base sur la recherche fondamentale
et permet par la suite les avanc\'ees techniques dans divers domaines
tels que la m\'edecine avec les techniques d'imagerie m\'edicales,
l'armement avec les armes et syst\`emes de d\'efense assist\'es au laser
ou encore l'industrie avec les nanotechnologies.
D'autres exemples peuvent \^etre donn\'es
dans les sciences de l'information et l'\'economie.

Nous voyons ici quelques challenges ce se propose de r\'esoudre 
notre branche de la physique dans les prochaines d\'ecennies.
On reconna\^it par exemple l'axe de recherche des coll\`egues qui 
travaillent sur les atomes froids et la condensation de Bose-Einstein.
Quant \`a nous, nous travaillons sur les derux derniers points, 
\`a savoir l'interaction laser atome avec intenses et super intenses 
\`a rayons x et UV ; et sur l'interaction avec les impulsions 
ultra-br\`eves qui pernmettront une pr\'ecision chirurgicale 
dans le contr\^ole des atomes et des mol\'ecules.

\item \textbf{Slide 3} \par
Les figures sur la gauche repr\'esentent la situation expr\'erimentale 
actuelle. 
On peut sur la premi\`ere figure qu'il est possible de produire
des inmpulsions laser aussi br\`eves que $5\,\rm fs$.
La deuxi\`eme figure nous montre essentiellement que les 
lasers \`a base de Titanium:Sapphire permettent pour ces impulsions 
tr\`es br\`eves de produire des champs d'intensit\'e sup\'erieure
\`a $100\,\rm TW/cm^2$.\par
Dans le m\^eme temps, la recherche th\'eorique 
conna\^it quelques difficult\'es dans la mesure o\`u 
on d\'enote un manque de mod\`eles th\'eoriques adapt\'es
pour al description de l'interaction rayonnement-mati\`ere 
en champ laser intense.
Par exemple, un tr\`es grand nombre de publications visant \`a d\'ecrire 
ces approximations sont encore effectu\'ees
dans l'espace des configurations.
Cet espace est inadapt\'e \`a la description de l'ionisation 
dans la mesure o\`u la particule \`a t\^ot fait d'aller \`a l'infini.
Afin de palier ce probl\`eme, on recquiert 
\`a plusieurs techniques d'approximation et de r\'eduction du 
domaine qui sont plus ou moins efficaces pour la r\'esolution
du probl\`eme mais ne s'appuient pas sur une base physique.
Toutefois, conform\'ement au principe d'incertitude, 
cette non localit\'e de la fonction d'onde dans l'espace des 
configuration entra\^ine une forte localisation de cette derni\`ere
dans l'esapce des moments, qui est adapt\'e pour la description 
de l'ionisation dans la mesure o\`u le monment de la particule 
y est fini.

\item \textbf{Slide 4} \par
Toutefois, cette sup\'eriorit\'e de l'espace des moments 
ne vient pas gratuitement dans la mesure o\`u 
l'\'equation de Schr\"odinger prend maintenant une forme 
int\'egro-diff\'erentielle et sa r\'esolution num\'erique 
recquiert maintenant la gestion du potentiel coulombien 
qui prend une forme non locale et comporte une  singularit\'e 
logarithmique.

Dans ce travail, nous d\'ecomposons le potentiel 
coulombien en une somme spectrale, facilement impl\'ementable 
du point de vue num\'erique et utilisons une technique 
de propagation qui rend la r\'esolution num\'erique 
de l'\'equation de Schr\"odinger viable.

Ultimement, ce mod\`ele permettra une description correcte
des exp\'eriences en champ laser au del\`a de l'approximation dipolaire 
qui ont lieu dans les infrastructures laser modernes.

\item \textbf{Slide 5}\par
Nous allons maintenant pr\'eciser le cadre th\'eorique dans
lequel nous travaillons avant de formellement pr\'esenter le mod\`ele 
qui repose au centre de ce travail.
Nous analyserons ensuite quelques r\'esultats issus de
la structure et de la dynamique avant de conclure.
\end{itemize}

\section{G\'en\'eralit\'es}

\begin{itemize}
\item \textbf{Slide 6}\par
Du fait de sa qualit\'e d'onde \'electromagn\'etique, 
la dynamique du champ laser est d\'ecrite par les quatre 
\'equations de Maxwell dans le vide.
Il est commode d'\'ecrire les champs \'electrique et magn\'etique 
\`a l'aide des potentiels scalaire et vecteur 
conform\'ement aux relations ci-apr\`es.
En utilisant le quadripotentiel relativiste, 
on \'ecrit alors l'\'equation d'onde \`a laquelle ob\'eit le rayonnement.
L'\'equation d'onde indique que les quadripotentiel 
d\'ecrivant les champs \'electrique et magtn\'etique
ne sont pas uniques. En effet, tout quadripotentiel $A_\mu^\prime$ 
li\'e \`a $A_\mu$ par la relation de jauge de premi\`ere esp\`ece
est abilit\'e \`a decrire le champ physique.
On parle alors de fixation de jauge.

Nous travaillons en champ fort, 
c'est \`a dire pour des champs laser d'intensit\'es sup\'erieures \`a 
$10^{12}\,\rm W/cm^2$.
Pour de telles intensit\'es, le potentiel vecteur 
peutr \^etre mis sous la forme suivante 
o\`u $f(t)$ est appel\'ee "l'enveloppe" et le terme en $\cos$ est appel\'e
"la porteuse".
On note que $\boldsymbol{\varepsilon}$ repr\'esente ici la 
polarisation de l'onde \'electromagn\'etique.

L'approximation dipolaire, repr\'esent\'ee par la condition 
$\boldsymbol{k}\cdot\boldsymbol{r}$, consiste \`a n\'egliger 
les dimensions de l'atomes devant la longueur d'onde de l'onde
\'electromagn\'etique. Son utilisation nous permet de 
mettre le rayonnement sous la forme suivante.

\item \textbf{Slide 7}\par
En utilisant cette description classique du rayonnement, 
on maintient une description quantique de la particule dont la dynamique 
est donn\'ee par son \'equation de Schr\"odinger.
On aprle alors de description semi-classique de
l'interaction rayonnement-mati\`ere.
L'hamiltonien $H$ dit "de couplage minimal" est alors la somme de 
l'hamiltonien  non-pertubr\'e $H_0$, 
du terme d'interaction paramagn\'etique 
et du terme d'interaction diamagn\'etique.

Le principe d'invariance de jauge nous permet alors de choisir les 
potentiels scalaire et vecteur tels que ${\rm div}\boldsymbol{A}=0$
et $\Phi=0$. On parle alors de jauge radiation.
Ce choix nous permet d'obtenir la forme usuelle de
l'hamiltonien en jauge vitesse donn\'ee par $H^v$.

D'autres hamiltioniens sont aptes \`a d\'ecrire 
le syst\`eme. En effet, tout hamiltonien 
$H^\prime$ li\'e \`a un hamiltonien $H$ 
d\'ecrivant le syst\`eme par la transformation de jauge 
de seconde esp\`ece, est lui-m\^eme abilit\'e \`a decrire le 
syst\`eme.

Par exemple, on obtient l'hamiltonien en jauge longeur 
g\'en\'eralement utilis\'e dans l'espace des configurations en 
posant $\chi=\boldsymbol{A}\cdot\boldsymbol{r}$.

\item \textbf{Slide 8}\par
Dans ces champs laser intenses, 
l'atome subit divers processus dont l'ionisation est le plus important. 
Keldysh introduisit un param\`etre d'adiabaticit\'e dont le calcul permet de 
pr\'eciser la dynamique de l'ionisation.
Ce terme d\'epend du potentiel d'ionisation $I_p$ 
de l'atome et du potentiel pond\'eromoteur $U_p$, lui-m\^eme fonction 
de l'intensit\'e du champ \'electrique et de la fr\'equence des
photons incidents.
Lorsque $\gamma \gg 1$, la dynamique est domin\'ee par 
l'ionisation au dessus du seuil, que nous d\'esignerons par 
processus ATI dans la suite.
Lorsque $\gamma \ll 1$, la dynamique est domin\'ee par 
un type d'ionisation dit "tunnel".
On note que pour des champs suffisamment forts,
le champ est suffisamment distordyu pour que la particule puisse 
effectivement \'echapper au puits de potentiel sans effet tunnel. 
On parle d'ionisation au dessus de  la barri\`ere.

On souligne \'egalement que le potentiel pond\'eromoteur $U_p$
permet de caract\'eriser les photo\'electrons qui 
sont qualifi\'es de "directs" lorsque leur \'energie est inf\'erieure
\`a $2U_p$ et "indirects", pour des \'energies comprises entre
$2U_p$ et $10U_p$.

\end{itemize}

\section{Nouvelle approche de r\'esolution dans l'espace des moments}

\begin{itemize}
\item \textbf{Slide 9}\par
Ayant plus explicit\'e l'essentiel des concepts pertinents 
\`a notre \'etude, nous voulons maintenant 
introduire le nouveau mod\`ele au c\oe ur de ce travail.

\item \textbf{Slide 10}\par
Les fonctions propres de la particule dans l'espace des configurations
sont fornm\'ees par un produit d'une fonction radiale 
et d'une fonction angulaire repr\'esent\'ee ici par l'harmonique sph\'erique.
La partie radiale, usuellement donn\'ee par un polyn\^ome de 
Laguerre peut \'egalement \^etre \'ecrite comme une fonction 
sturmienne $S_{n\ell}^\kappa$.
Ces fonctions ob\'eissent aux conditions aux limites usuelles.

On note que les fonctions sturmiennes sont explicitement donn\'ees 
par la relation suivante o\`u ${}_1F_1$ est la fonction hyperg\'eom\'etrique 
confluente et $\widetilde{N}_{n\ell}^\kappa$ est un facteur de 
normalisation.

Les fonctions propres de d\'egr\'es sup\'erieurs  
peuvent \^etre construites en utilisant la relation de r\'ecurrence suivante
baqs\'ee sur celle des fonctions hyperg\'eom\'etriques confluentes.

La figure repr\'esente la variation des fonctions sturmiennes 
en fonction du param\`etre non lin\'eaire $\kappa$.
On peut voir que l'augmentation de $\kappa$ r\\'eduit l'extension 
spatiale de la fonction sturmienne qui varie \`a priori de 0 \`a l'infini, 
pendant que son amplitude semble bien log\'ee entre 0 et 1.

\item \textbf{Slide 11}\par
Ce constat va nous aider \`a cerner l'int\'er\^et de l'espace des moments 
dont les fonctions propres s'expriment aussi comme 
le produit d'une fonction radiale et d'uen harmonique sph\'erique modifi\'ee.
La relation liant les parties radiales dans les espaces des configurations 
et des moments est donn\'ee par la relation suivante.  

La fonction sturmienne radiale est explicitement donn\'ee par cette 
relation, o\`u $C_{\bar n}$ est le polyn\^ome de Gegenbauer
et $\mathcal{N}_{n\ell}^\kappa$ est un facteur de normalisation.

Les fonctions radiales d'ordre sup\'erieures sont b\^aties sur la relation 
de r\'ecurrence suivante bas\'ee sur celle des polyn\^omes de Gegenbauer.

Sur la figure est repr\'esent\'ee l'effet sur les fonctions 
sturmiennes de la variation du param\`etre non lin\'eaire $\kappa$.
On peut voir que $\kappa$ a ici un effet sur l'amplitude 
de la fonction d'onde, tandis que son extension spatiale
reste localis\'ee entre -1 et 1.
Cette r\'eduction du domaine auparavant de 0 \`a l'infini
\`a l'intervalle $\left[-1;+1\right]$ ext exactement ce qui rend l'espace des 
moments attractif sur la plan num\'erique.

\item \textbf{Slide 12}\par
La particule est d\'ecrite dans l'espace des configurations 
par la fonctions $\varphi$ satisfaisant \`a l'\'equation de Schr\"odinger 
et qui s'\'ecrit comme une somme des fonction de base de l'espace.

La formulation dans l'espace des moments peut \^etre obtenue 
par une transformation de Fourier de toute l'\'equation pr\'ec\'edente 
et s'\'ecrit comme suit.

On note l'apparition de l'int\'egrale dans l'\'equation.
Ainsi, la solution de l'\'equation en un point $\boldsymbol{p}$
de l'espace d\'epend de la solution en tous les autres points 
$\boldsymbol{p}^\prime$ du domaine. 

Une attention particuli\`ere  est port\'ee \`a la nouvelle 
forme prise par le potentiel coulombien avec la singularit\'e
logarithmique au d\'enominateur.

Nous avons pu montrer que cette singularit\'e logarithmique 
peut se mettre sous la forme d'une somme spectrale ici donn\'ee, 
ce qui nous permet de r\'e\'ecrire l'\'equation de Schr\"odinger 
comme suit.

\item \textbf{Slide 13}\par
L'int\'er\^et de cette derni\`ere \'equation se manifeste lorsque 
l'on passe \`a la repr\'esentation matricielle de l'hamiltonien.
Le probl\`eme aux valeurs propres g\'en\'eralis\'e qui en d\'ecoule
s'\'ecrit donc comme suit, o\`u $S$ est la matrice de recouvrement
et l'hamiltonien non perturb\'e $H^{(0)}$ est la somme de 
l'\'energie cin\'etique $\mathcal{T}$ 
et de l'\'energie potentielle coulombienne $\mathcal{V}$.

Le d\'etail de ces termes est donn\'e par ces trois \'equations.
On voit \`a partir de ces derni\`eres que les matrices de recouvrement 
et de l'hamiltonien non perturb\'e sontb tridiagonales par bloc.

\item \textbf{Slide 14}\par
L'hamiltonien total $H$ \'etant la somme de l'hamiltonien non perturb\'e 
$H^{(0)}$ et de l'hamiltonien d'interaction $H^{(1)}$, 
on calcule ce dernier donn\'e ici dans la jauge vitesse et 
pour un champ laser polaris\'e de fa\c con rectiligne.
On obtient les termes suivants et la structure de la matrice se d\'eduit
encore une fois des symboles de Kronecker. 

Ainsi, en ce qui concerne l'hamiltonien total $H$, on voit que seul les blocs 
$\Delta\ell=0,\pm 1$ sont occup\'es et que le nombre quantique magn\'etique 
$m$ se conserve au cours de l'approximation. 
Ces r\'esultats sont bien ceux pr\'edits par la th\'eorie
pour l'utilisation  d'une polarisation lin\'eaire 
dans l'approximation dipolaire.


\item \textbf{Slide 15}\par
L'hamiltonien total ainsi construit dans la base sturmienne, 
la propagation dans le temps du vecteur solution se fait \`a l'aide 
d'un sch\'ema d'int\'egration num\'erique.

Par soucis de performance, nous proposons une technique de propagation bas\'ee 
sur le sch\'ema publi\'e par Nurhuda et Faisal dans leur article de 1999.
En utilisant les formules de changement de base de l'alg\`ebre lin\'eaire,
on peut \'ecrire \`a tout instant l'hamiltonien et le vecteur d'onde solution
la base atomique \`a l'aide des relations qui suivent.

On \'ecrit ainsi la solution \`a un instant $t$ ult\'erieur
en utilisant la m\'ethode de Crank-Nicholson.
La matrice au d\'enominateur doit \^etre invers\'ee \`a chaque pas de 
propagation, ce qui est num\''eriquenment tr\`es co\^teux.
On \'ecrit donc cette matrice comme la somme d'une matrice diagonale 
$O_D$ et d'une matrice non diagonale $O_{ND}$ et on fait un d\'eveloppement 
en s\'erie de Taylor pour \'ecrire la formule qui vient.

Comme on peut le voir, dans cette derni\`ere formule, 
le terme qui n\'ecessite maintenant une inversion \`a chaque pas de propagation  
est la matrice diagonale $O_D$ qui est mod\'elis\'e e par un vecteur.

Le gain num\'erique occasionn\'e par cette technique par rapport \`a un sch\'ema 
Runge-Kutta d'ordre 4 classique et de l'ordre de 50\%.

\item \textbf{Slide 16}\par
La vecteur d'onde solution obtenu \`a la fin de la propagation dans la base atomique, 
on utilise les m\^emes formules de changement de base pour \'ecrire
la fonction d'onde dans la base sturmienne.

Diverses observables telles que les populations des \'etats li\'es, 
les distributions en moment, en \'energie 
et en impulsions transversales et longitudinales peuvent ensuite \^etre \'evalu\'ees
\`a la fin de l'interaction ou au cours de celle-ci 
par les diverses formules donn\'ees sur ce slide.

\section{R\'esultats}

\item \textbf{Slide 17}\par
Nous allons maintenant \'eprouver ce nouveau mod\`ele en 
analysant les r\'esultats du point de vue de la structure, 
mais aussi de la dynamique dans les hautes et basses fr\'equences.

\item \textbf{Slide 18}\par
Nous avons construit l'hamiltonien de la particule 
dans une repr\'esentation alternative qui est la base sturmienne dans 
l'espace des moments. Nous voulons donc tout d'abord nous assurer 
que l'hamiltonien que nous calculons repr\'esente bien notre syst\`eme, 
c'est-\`a-dire, l'atome d'hydrog\`ene.

Pour s'esn assurer, nous allons comparer les valeurs propres 
obtenues par la diagonalisation de notre hamiltonien aux valeurs 
th\'eroriques exactes.
Les r\'esultats sont consign\'es dans le premier tableau.

Comme on peut le voir, la pr\'ecision diminue avec l'ordre $n$
du niveau consid\'er\'e.
Pour les param\`etres choisis, la pr\'ecision est mauvaise \`a partir 
du 70e niveau.
De fa\c con g\'en\'erale, 
la taille de base choisie influe sur le nombre d'\'etats li\'es 
d\'ecrits et 60\% des ces \'etats li\'es d\'ecrits le sont avec une pr\'ecision 
de plus de 10 chiffres. 

%On note que les param\`etres de base ici utilis\'es 
%permettent de d\'ecrire 83 \'etats li\'es.

Alternativement, la qualit\'e de la description dus syst\`eme 
peut \^etre \'evalu\'ee en comparant les fonctiosnm propres obtenues num\'eriquement 
aux valeurs de r\'ef\'erence.
Les r\'esultats sont consign\'es dans le tableau ci-bas.
On peut voir que le recouvrement g\'en\'eral est sur 11 chiffres, 
ce qui augure d'une description tr\`es pr\'ecise du syst\`eme.

\item \textbf{Slide 19}\par
Finalement, on peut \'evaluer la structure en comparant les 
forces d'oscillateur calcul\'ees \`a l'aide des valeurs et vecteurs propres 
de notre syst\`eme aux valeurs de r\'ef\'erences.

Le tableau \`a double entr\'ees suivant comprare les valeurs de 
r\'ef\'erence (pr\'esent\'ees sur la premi\`ere entr\'ee)
aux valeurs calcul\'ees par la r\'esolution num\'erique 
de l'\'equation de Schr\"odinger (deuxi\`eme entr\'ee).
On voit que l'on reproduit
la struture de l'hydrog\`ene dans les limites de la double pr\'ecision 
sur pr\`es de 50 niveaux, ce qui est une tr\`es bonne description de la structure.

Accessoirement, on peut \'evaluer les polarisabilit\'es calcul\'ees 
\`a partir de notre mod\`ele \`a celles donn\'ees par la formule th\'eorique.
Le tableau ci-bas pr\'esente ces r\'esultats pour de l'hydrog\`ene atomique 
et l'ion h\'elium dans son \'etat initial $2s$.
On peut voir que l'on reproduit parfaitement les valeurs th\'eoriques,

\item \textbf{Slide 20}\par
Ces quelques tests de validation de la structure auyant \'et\'e effectu\'es, 
nous pouvosn maintenant \'etudier les premiers r\'esultats de
la dynamique donn\'es par notre mod\`ele en r\'egime haute fr\'equence.

Sur les deux figures suivantes, nous comparons les r\'esultats 
issus de notre mod\`ele \`a ceux obtenus par la M\'ethode des 
Coordonn\'ees D\'ependante du Temps de Hamido \textit{et al.}
On note que cette m\'ethode utilise l'espace des configurations.

\item \textbf{Slide 21}\par
Dans la m\^eme id\'ee, nous comparons maintenant les r\'esultats 
issus de ces deux mod\`eles \`a celui des potentiels s\'eparables 
de Tetchou \textit{et al.}
On peut voir que les r\'esultats issus de cette derni\`ere m\'ethode
diff\`erent significativement de notre mod\`ele et de la MCDT.

La raison est que ce mod\`ele ne consid\`ere que deux niveaux 
du potentiel coulombien, et de ce fait est comparable \`a l'approximation du champ 
fort qui n'en consid\`ere qu'un seul.

Cesn r\'esultats illustrent donc bien le besoin d'int\'egrer 
de fa\c con satisfaisante le potentiel coulombien dans la dynamique.

\item \textbf{Slide 22}\par
Nous voulons maintenant analyser un peu plus en profondeur 
les r\'esultats de la dynamique issus de notre mod\`ele 
dans le proche infrarouge pour des champs laser de diff\'erentes dur\'ees 
et pour des impulsions de diff\'erentes formes.

Les signaux que nous s\'electionnons sont de trois types et pour 
chacun d'entre eux, nous faisons varier un nombre r\'eduit de param\`etres 
et analysons les effets sur la physique de l'ionisation.

Il est important de noter que chacune des courbes que nous
pr\'esentons ici les r\'esultats de plus de 22 runs et 
que l'obtention de la plupart de ces courbes n\'ecessite pr\`es de deux semaines 
de calculs sur les centres calculs sp\'ecialis\'es.

\item \textbf{Slide 23}\par
Nous commen\c cons donc par l'analyse des populations des
\'etats li\'es et des probabilit\'es d'ionisation en fonction du temps 
en fonction de l'intensit\'e.
Sur la figure de gauche, les courbes noire et rouge repr\'esentent 
respectivement les r\'egimes multiphotonique et interm\'ediaire.
On peut voir qu'\`a la fin de l'interaction, l'\'etat fondamental 
$1s$ ne se d\'epeuple pas.
Pour le r\'egime tunnel repr\'esent\'e par la courbe bleue en revanche, 
on peut voir qu'une part significative (20\%) des particules est ionis\'ee.
La deuxi\`eme figure illustre la variation des ces populations 
pour un cas d'ionisation tunnel en noir, et d'ionisation au dessus du seuil
(ren rouge). 
On voit que dans ce dernier cas, une part encore plus importante de 
particule est ionis\'ee \`a la fin de l'interaction.

Ces r\'esultats sont attendus dans la mesure o\`u il existe une intensit\'e 
de saturation $I_s$ pour laquelle la totalit\'e des particules
est ionis\'ee \`a la fin de l'interaction.

\item \textbf{Slide 24}\par
Dans la mesure o\`u nous n'avons fait de comparaisons avec d'autres 
auteurs qu'\`a haute fr\'equence, on peut se demander ce que donneraient 
de telles comparaisions pour l'infrarouge proche.

Les tableaux repr\'esentent la comparaison de nos r\'esultats 
avec ceux donn\'es par deux diff\'erentes techniques de Chen \textit{et al.}
Le recouvrement de nos r\'esultats avec les leurs est de l'ordre du dixi\`eme.
On voit sur le pre4mier tableau que pour les impuslions br\`eves, 
nos r\'esultats sont plus proches 
de ce qu'ils obtiennent avec leur m\'ethode spectrale, tandis que pour les 
impulsiosn plus longues, sur le deuxi\`eme tableau, l'accord est meilleur
avec leur m\'ethode des grilles.

C'est l'occasion de souligner la pr\'ecision avec laquelle nous travaillons 
(sur 16 chiffres) par rapport \`a la plupart des autres auteurs 
qui donnent des r\'esultats au milli\`eme tout au plus.

Outre cela, la figure donne les populations des couches $n$ 
et des sous-couches $\ell$ pour diff\'erentes dur\'ees de l'impulsion 
br\`eves en haut et longues en bas, et pour diff\'erentes phases.
On peut y voir que les effets de phases ne s'expriment que pour les impulsions 
br\`eves dans la mesure o\`u les populations finales sont phase-d\'ependantes.
Pour les impulsions longues en revanche, on voit que les trois distributions 
relatives aux trois phases convergent.

Ceci nous emm\`ene \`a un r\'esultat bien connu de l'interaction 
rayonnement-mati\`ere : pas d'effets de phase pour les impulsions longues.

\item \textbf{Slide 25}\par
Il nous est \'egalement possible d'analyser le d\'etail de chacun 
des populations de chacun des \'etats $n\ell$ \`a la fin de 
l'interaction.

La premi\`ere figure \`a gauche donne les populations finales 
pour deux diff\'erentes formes d'impulsions br\`eves IL1 sur la premi\`ere colonne
et IL2 sur la deuxi\`eme, pour trois diff\'erentes phases
$\phi=0,\pi/4,\ \mbox{et}\ \pi/2$.
On voit tr\`es clairement la d\'epedance des populations finales 
aux phases consid\'er\'ees.

Sur la deuxi\`eme figure, nous \'etudions les populations 
des niveaux et sous-niveaux pour les diff\'erentes r\'egimes multiphotonique,
interm\'ediaire et tunnel. On peut y voir que les amplitudes 
des distributions augmentent avec l'intensit\'es et que les sous \'etats 
$n,\ell>8$ ne se peuplent pas malgr\'e l'augmentation de l'intensit\'e. 

Sur la troisi\`eme figure, on peut voir les distributions 
pour des impulsions tr\`es longues de 20 cyles optiques pour ${\rm H}(1s)$,
${\rm He^+}(2s)$ et ${\rm He^+}(2p)$.
On note que le comportement de ${\rm He^+}(2p)$ est tr\`es semblable 
\`a celui de ${\rm H}(1s)$.
 
\item \textbf{Slide 26}\par
De fa\c con alternative, on peut analyser la dynamique par le biais de 
d'autres observables. 
La dynamique de l'ionisation est tr\`es bien repr\'esent\'ee 
par l'\'etude des densit\'e d'\'energie en fonction du temps 
comme montr\'e sur la figure suivante.
Les courbes bleues rep\'esentent le potentiel vecteur et les r\'egions
en jaune repr\'esentent celles dens\'ement peupl\'ees en \'electrons.
On peut voir que pour ces figures repr\'esentant les trois r\'egimes, 
l'ionisation survient de plus en plus t\^ot avec l'augmentation de l'intensit\'e.
En effet, on constate l'apparition des raies \`a partir du 4e, 3e et 2e cycle optique,
en allant de gauche \`a droite. 
On remarque aussi les raies verticales noires fixes par rapport \`a l'axe $p_z$
qui donnent en fait les positions des pics ATI.

\item \textbf{Slide 27}\par
L'\'etude des densit\'es d'\'energie \`a la fin de l'interaction 
nous donne les informations usuelles sur les effets de phase 
pour les impulsions br\`eves.
La premi\`ere colonne repr\'esente ces densit\'es d'\'energie 
relatives \`a trois diff\'erentes phases pour une impulsion de type IL1.
On peut voir que ind\'ependamment de ces phases, la distribution vers 
l'arri\`ere domine ; tandis que pour uen impulsion de type 
IL2, la distribution arri\`ere domine pour $\phi=0$ et $\pi/4$, 
mais pas pour $\phi=\pi/2$. 

On peut voir sur la figure suivante l'\'evolution des
effets de phase avec l'augmentation de la dur\'ee de l'impulsion.
On voit clairement la disparition de l'asym\'etrie gauche/droite.


\item \textbf{Slide 28}\par
Nous voulons maintenant \'etudier l'interaction rayonnement-mati\`ere 
par le biais des spectres d'\'energie.
La figure de gauche nous montre l'\'evolution de ces spectres 
avec la dur\'ee d'impulsion. 
On voit tr\`es clairement que les probabilit\'es d'ionisation augmentent 
avec la dur\'ee de l'interaction.
Toutefois, on remarque que la r\'egion concernant les \'electrons directs
reste inaffect\'ee par la dur\'ee du pulse.

Ceci contraste avec la figure \`a droite qui montre 
l'\'evolution des spectres d'\'energie avec l'augmentation de l'intensit\'e 
du laser. 
On constate que cette fois ci, l'augmentation de la probabilit\'e 
d'ionisation est uniforme.


\item \textbf{Slide 29}\par
Ces figures repr\'esentent les spectres en \'energie de
diff\'erentes esp\`eces chimiques pour de tr\`es longues impulsions de 
20 cycles optiques.

Sur les bords, on peut voir les spectres 
relatifs \`a l'hydrog\`ene initialement dans son \'etat fondamental 
\`a l'extr\^eme gauche
et celui de l'\'helium sur son \'etat initial $2p$ \`a l'extr\^eme droite.
Ces spectres sont similaires par leur forme qui, 
ont le montre, est une signature de l'hydrog\`ene.
On observe le contrast avec les spectres au centre qui repr\'esentent 
qui repr\'esentent l'helium dans son \'etat initial $2s$.
La raison pour laquelle l'h\'elium $2p$ pr\'esente des similitudes 
qualitatives avec l'hydrog\`ene est mal comprise.

Outre cela, les deuxi\`eme et troisi\`eme lignes
repr\'esentent les \'emissions vers l'avant et vers l'arri\`ere pour 
ces longs pulses.
Ces spectres \'etant presque indistingables, 
ils t\'emoignent de l'absence des effets de phase pour ce type 
d'impulsion.


\item \textbf{Slide 30}\par
Un tr\`es grand nombre de conclusions tir\'ees de l'analyse 
des pr\'ec\'edentes observables est visible sur les distributions
en moments que l'on voit sur ces figures.

Par exemple, l'asym\'etrie gauche/droite dont il a d\'ej\`a \'et\'e 
question auparavant est bien visible sur la figure de gauche.
On voit bien la domination des $p_z$ n\'egatifs 
ind\'enpdamment de la phase pour la premi\`ere colonne 
et le redressement progressifs vers les $p_z$
positifs pour la deuxi\`eme.

Sur la deuxi\`emje figure, on voit l'\'evolution de
distributions en moments avec la dur\'ee de l'impulsion.
On voit clairement la sym\'etrisation de la distribution avec 
la longueur de l'impulsion.

Sur la troisi\`eme figure, 
on voit ces distributions pour les trois r\'egimes d'ionisation.
On peut voir la diff\'erence des probabilit\'es d'ionisation 
\`a travers les amplitudes sur les l\'egendes \`a droite.

On note aussi que les raies circulaires sont une manifestation 
des spectres ATI que nous avons rencontr\'es pr\'ec\'edement.


\item \textbf{Slide 31}\par
Finalement on peut observer l'int\'egration des distributions en 
moments sur les impulsions transversale et longitudinale.

Sur cette derni\`ere figure, on peut voir l'immense domination
de la r\'egion $p_z < 0$ ind\'ependamment de la phase, et l'\'evolution 
progressive vers les $p_z > 0$ sur la deuxi\`eme figure.

Sur ces autres figures, on peut voir l'\'evolution de 
cette m\^eme observable pour les impulsions longues avec l'intensit\'e.
On y voit l'habituelle augmentation des probabilit\'es 
d'ionisation ainsi qu'une sorte de l\'eg\`ere sym\'etrisation des distributions 
avec l'intensit\'e.

\section{Conclusion}

\item \textbf{Slide 32}\par
Nous avons analys\'e l'interaction
pour diff\'erents types d'impulsions br\`eves et longues et dans 
diff\'erents r\'egimes d'intensit\'es.
Ces \'etudes sont n\'ecessaires pour 
saisir l'effet du rayonnement sur les la mati\`ere pour les impuslions 
ultra-br\`eves et ultra-intenses qui sont appel\'ees \`a \^etre utilis\'ees 
en m\'edecine et dans l'industrie.

\item \textbf{Slide 33}\par
Nous avons men\'e cette \'etude dans l'espace des moments 
du fait de sa sup\'eriorit\'e th\'eorique pour la description de 
l'ionisation par rapport \`a l'espace des configurations.

Cet espace posant des probl\`emes connus, nous avons mis au point 
un nouvelle approche qui consiste \`a \'ecrire le potentiel coulombien 
non local comme une simple somme spectrale.
La r\'esolution de l'\'equation de Schr\"odinger s'en trouve alors 
grandement simplifi\'ee et acc\'el\'er\'ee sur la plan num\'erique.
De plus, nous avons mis au point un sch\'ema de 
propagation specialement adapt\'e \`a notre probl\`eme exploitant les
propri\'et\'es de la base atomique pour obtenir des 
performances qui rendent l'\'etude dans la r\'egions des XUV envisageables.

Nous nous omme ensuite assur\'es d'une coh\'erence interne de
nos r\'esultats en \'evaluant nos observables par diff\'erentes techniques 
avant de les confronter aux autres r\'esultats pr\'esenmts dans la litt\'erarture.
Nous obtenons d'excellents r\'esultats jusqu'\`a pr\'esent.

\item \textbf{Slide 34}\par
Dans l'avenir, 
on cherchera \`a appliquer ce mod\`ele dans les r\'egimes 
o\`u il fera une diff\'erence majeure, c'est \`a dire dans les r\'egimes
XUV de super hautes intensit\'es, au del\`a de l'approximation dipolaire.  

Nous avons pu montrer que notre d\'eveloppement est r\'ealisable 
dans la base dite des B-Splines, ce qui traduite une certaine 
versatilit\'e de notre approche.

Des \'etudes sont r\'ealis\'ees sur la construction 
d'un nouveau mod\`ele \`a potentiels s\'eparables b\^ati sur notre approche, 
qui nous permettra \'eventuellement d'\'etudier la contribution de chacun des niveaux 
d'\'energie comme l'ont fait avons nous Tertchou et ses partenaires.

Et pour finir, notre mod\`ele \'etant valables pour les hydrog\'eno\"ides 
en g\'en\'eral, nous sommes en train de l'appliquer \`a diff\'erents
syst\`emes \`a un \'electron actif.


\item \textbf{Slide 35}\par
Ce travail a fait l'objet de quelques publication, l'une parue depuis 2016,
les autres ayant \'et\'e soumises cette ann\'ee 
dans des journaux internationaux de physique atomique.
atomique. 

\item \textbf{Slide 36}\par
je gfinirai donc ce talk en remerciant mon Universit\'e h\^ote, 
l'Universit\'e des Sciences et Technique de Masuku, 
en particulier la Facult\'e des Sciences et l'Ecole Polytechnique 
pour leur soutien \`a mon travail.

Je ne saurais finir sans remercier le Centre de Physique Atomique, 
Mol\'eculaire et d'Optique Quantique de l"universit\'e de Douala dans lequel 
j'ai effectu\'e plusieurs stages de recherches et de perfectionnement.

Enfin, je remercierai l'Agence Nationale des Bourses du Gabon 
qui a certainement \'et\'e ma seule source de financement le long de ces recherches.

Je vous remercie.
\end{itemize}

\end{document}
